\section{Introduction}
\label{sec:intro}
Hybrid retrieval for retrieval-augmented generation (RAG) commonly combines the Top-$L$ results of dense and sparse retrievers using reciprocal rank fusion (RRF)~\cite{elastic,rrf}. A fixed window limits reads, but also removes contributions beyond that window. When a document appears in one channel but not in the other channel's window, fusion assigns zero to the latter contribution. Yet absence from the window only means that the rank has not been observed: its contribution is unknown. Replacing it with zero can change the output order even when the observed candidates already contain every document in the complete-fusion Top-$K$.

Diverse, open-ended requests and continuing corpus updates change channel rankings and their overlap, making a fixed $L$ difficult to transfer across queries and corpus snapshots. Exact fusion avoids a fixed window by bounding unread contributions and reading until the first $K$ documents and their order are determined~\cite{nra}. EAHR implements this approach for hybrid retrieval~\cite{eahr}. However, the reads needed to complete a fixed Top-$K$ can still vary sharply across queries. Resolving a few contested positions can require reading far down the channel rankings.

Resource-constrained requests therefore need a direct limit on reads. Specifying $K$ determines how many results must be completed, but does not bound the work required to complete them. We instead take a read budget as input and return the entire certified prefix at that budget. Execution determines the output length, while its order remains identical to complete fusion.

We study this budgeted exact-prefix formulation through DiBud (Direct Budgeting). DiBud bounds unread contributions, prioritizes the channel blocking the next position, and emits results as they become certain. It combines selective reading from SNRA~\cite{snra} with incremental exact output from LARA-IN~\cite{lara}. Replays on five query sets compare the certified yield of different schedules under equal budgets; a further replay across five corpus snapshots examines delivery as the corpus changes.
